% !TEX root = ../main.tex
\chapter*{概説}\label{gaisetu}\addcontentsline{toc}{chapter}{概説}\thispagestyle{plain}\setcounter{chapter}{1}

\section{水納島方言}
南琉球宮古語\ruby[g]{水納}{みんな}\ruby{島}{じま}方言（自称〈みんなふつ\heiban{}〉）は沖縄県宮古郡多良間村水納島で伝統的に話されていたことばである。1962年より集団移住の政策が取られたため、島民のほとんどが宮古島の高野集落に移った。現在も水納島方言の多くの話し手たちが高野や宮古島の他の集落に住んでいる。しかし、島民たちが集落を離れて60年間経っている今は水納島の伝統的なことばが全く継承されておらず、その上に伝統的な方言の最後の継承者が少なくなっている。

%水納島は多良間島の北に浮かぶ小島で、宮古島と石垣島の間に位置している。宮古島から約67キロメートル、石垣島から約35キロメートル離れている。水納島


水納島方言は系統的に宮古語に属しており、その中で特に多良間島方言（\ruby[g]{仲筋}{なかすじ}方言、\ruby[g]{塩川}{しおかわ}方言）に近いことが分かっている（図1）\citep{lawrence03:taramakeito,pellard09:phdthesis,celik20:phdthesis,celik20:minna}。

\begin{figure}[ht!]\centering
\begin{tikzpicture}
\tikzset{level distance=0.8cm}
\tikzset{sibling distance=3cm}
%\coordinate
\node{日琉諸語}
  child { node {琉球諸語} 
    child { node{南琉球} 
      child { node{八重山語} }
      child { node{宮古語}
       child { node{多良間諸方言} 
        child { node{\underline{水納島方言}} }
        child { node{多良間島方言} }
      }
       child { node{共通宮古語} }
      } 
    }
    child { node{北琉球} } 
  }
  child { node{日本語}
  }
;
\end{tikzpicture}
\caption{水納島方言の系統}\label{fg:keito}
\end{figure}


多良間諸方言を構成する水納島方言と多良間島方言は音韻論が大きく異なるものの、それ以外の点においてよく似ており、相互理解が完全に成り立っている。明治と大正生まれの水納島方言の報告\citep{zenshuraku:minna}を見ると、明治期と大正期との間に大きな音変化があったと推測される。つまり、〈イ゜〉の母音が〈さ〉〈ざ〉〈つぁ〉の行を除いて〈い〉の母音と完全に合流している。明治生まれの話者が〈ぽーキ゜〉「箒」と発音していたところ、現在は〈ぽーき〉と発音する。一方、興味深いことに水納島方言のアクセント体系の方が多良間島方言に比べ保守的な側面が多い。


% 話者情報

\section{発音}

水納島方言の子音体系を表\ref{tb:shiin}に示す。注意点として、〈さ〉行、〈ざ〉行と〈つぁ〉行の発音の揺れがある。特に〈さ〉行と〈ざ〉行は非口蓋化の発音（さ、すぅ、せ、そ／ざ、ずぅ、ぜ、ぞ）と口蓋化した発音（しゃ、しゅ、しぇ、しょ／じゃ、じゅ、じぇ、じょ）のとで激しく揺れている。〈つぁ〉行も揺れる場合があるが、〈ちゃー〉「茶」のように揺れない語もある。

\begin{table}[ht]\caption{子音体系}\label{tb:shiin}
%\resizebox{\textwidth}{!}{    
\centering
\def\arraystretch{0.95}
     \begin{tabular}{llllll}
     \toprule
        & 両唇・唇歯音 & 歯茎音 & 軟口蓋音 & 声門音 \\
        \midrule
        破裂音 & ぱ \ip{[pa]}  & た \ip{[ta]}  & か \ip{[ka]} &   \\
        & ば \ip{[ba]} & だ\ip{[da]} & が \ip{[ga]} \\[2pt]
        鼻音 & ま \ip{[ma]} & な \ip{[na]} & &  \\[2pt]
        流音 & & ら \ip{[ra\ntilde{}ɾa]} & & \\[2pt]
        摩擦音 & ふぁ \ip{[fa]} & さ \ip{[sa\ntilde{}ɕa]}& & （は \ip{[ha]}） \\
        & ヴぁ \ip{[va\ntilde{}ʋa]} & & \\[2pt]
        破擦音 & & つぁ \ip{[tsa\ntilde{}tɕa]} & & \\
        & & ざ \ip{[dza\ntilde{}dʑa]} & & \\[2pt]
        接近音 & & や \ip{[ja]} & & \\
        %\midrule
        %二重子音 & & & & & & \\
     \bottomrule
 \end{tabular}
% }
\end{table}

もう1つの注意点として、〈ん〉と〈ム〉の区別が曖昧になっていることが挙げられる。つまり、〈ム〉がよく〈ん〉と発音される。例えば、過去接辞の〈たム〉が〈たん〉ともよく発音されるわけである。

水納島の母音体系を表\ref{tb:boin}に示す。注意点として〈イ\m{}〉の母音がある。この母音は〈さ〉行、〈ざ〉行、〈つぁ〉行の子音としか共起しない。つまり、〈す〉〈ず〉〈つ〉以外の音節がない。また、〈い〉の発音と揺れている。〈うす〉「牛」は「うし」、〈みず〉「水」は〈みじ〉、〈みつ〉「道」は〈みち〉とも発音できるわけである。ただし、元の〈い〉の母音は基本的に揺れず、〈ちび〉「尻」は〈ちび〉とだけ発音される。

\begin{table}[ht]\caption{母音体系}\label{tb:boin}
%\resizebox{\textwidth}{!}{    
\centering
\def\arraystretch{0.95}
     \begin{tabular}{lllllll}
     \toprule
        狭 & い \ip{[i]} & いー \ip{[iː]} & イ\m{} \ip{[ʉ]} & イ\m{}ー \ip{[ʉː]}  & う \ip{[u]} & うー \ip{[uː]}  \\[2pt]
        半狭 & え \ip{[e]} & えー \ip{[eː]} & & & & おー \ip{[oː]} \\[2pt]
        広 & & & & & あ \ip{[a]} & あー \ip{[aː]} \\[2pt]
        \bottomrule
 \end{tabular}
% }
\end{table}


\FloatBarrier

\section{アクセント}

水納島方言は主に3つのアクセント型（a型、b型、c型）を区別している\citep{celik20:minna}。これらのアクセント型はピッチ（声の高さ）の局所的な変動（下降または上昇）の有無と位置によって対立している。a型はピッチの変動を伴わないのに対して、b型は遅めのピッチの変動、c型は早めのピッチの変動を伴う。それに加えて、数語にしか観察されない特殊型も存在する。この特殊型は2つのピッチの変動（下降と上昇）を伴う。それぞれのアクセント型の例を\Next[]に示す（ピッチ変動のないことを「\heiban{}」、ピッチの下降を「\kako{}」、ピッチの上昇を「\josho{}」の記号で表す）。

\ex. 
    \makebox[4em][l]{a型：}\makebox[11em][l]{ぶとぅまい\heiban{}　ぶらーん}「夫もいない」\\
    \makebox[4em][l]{b型：}\makebox[11em][l]{みどぅムま\kako{}い　ぶらーん}「女性もいない」\\
    \makebox[4em][l]{c型：}\makebox[11em][l]{う\kako{}やまい　ぶらーん}「お父さんもいない」\\
    \makebox[4em][l]{特殊型：}\makebox[11em][l]{ヴぇー\kako{}い\josho{}まい　あらん}「これでもない」


しかし、多良間方言と同様に、各アクセント型の実現が環境によって異なる。上の\Last[]の例は発話のはじめという環境に該当しているが、発話の頭にc型の〈ん\kako{}め〉「もう」を入れてみると、\Next[]のように各アクセント型の実現が変わってくる（特殊型は省略する）。つまり、a型が低く平らに発音されるのに対して、b型とc型はそれぞれ遅めと早めのピッチの上昇で実現している。

\ex. 
    a型：\makebox[15em][l]{ん\kako{}め　ぶとぅまい\heiban{}　ぶらーん}「もう夫もいない」\\
    b型：\makebox[15em][l]{ん\kako{}め　みどぅムまい　\josho{}ぶらーん}「もう女性もいない」\\
    c型：\makebox[15em][l]{ん\kako{}め　うや\josho{}まい　ぶらーん}「もうお父さんもいない」


\LLast[]におけるアクセント型の実現を「下降調」、\Last[]におけるものを「上昇調」と呼ぶ。各アクセント型の実現がこのように環境によって逆転するという特徴は日琉諸語の中でもたいへん珍しく、他に報告がないようである。

\section{動詞の活用クラス}

動詞は2種類の規則的な活用類、すなわち\rom{1}類と\rom{2}類に分かれる。\rom{1}類に属する動詞は「書く」「読む」のように日本語の五段動詞に対応しており、命令形が〈い〉に終わり、否定形において〈あ〉が現れる。\rom{2}類に属する動詞は「見る」「下りる」のように一段動詞に対応しており、命令形が〈る〉に終わり、否定形において〈あ〉が現れない。それに加えて、混合型の動詞〈すに〉「死ぬ」と不規則な活用を示す変則型の動詞〈しー〉「する」と〈きー〉「来る」の動詞がある。表\ref{tb:katuyorui}にそれぞれの活用類の代表例を示す。

\begin{table}[ht]\caption{活用類の代表例}\label{tb:katuyorui}
\resizebox{\textwidth}{!}{    
\centering
\def\arraystretch{0.9}
     \begin{tabular}{lllll}
     \toprule
        活用類 & 動詞 & 基本形  & 命令形 & 否定形　\\
        \midrule
        \rom{1} &「書く」& かき  & かき & かかん \\
         & 「漕ぐ」&くぎ& くぎ & くがん \\
        & 「飛ぶ」 &とぅび& とぅび& とぅばん\\
        & 「読む」 &ゆム& ゆみ & ゆまん\\
        & 「待つ」 &まつ／まてぃ& まてぃ & またん\\
        & 「成す」 &なす & なし & なしゃん\\
        & 「成る」 &ない & なり & ならん \\
        & 「実る」 &みーり & みーり & みーらん \\
        & 「降る」 &っふぃ & っふぃ& っふぁん\\
        & 「売る」&っヴぃ & っヴぃ& っヴぁん\\
        & 「移る」 &うっつ & うっち & うっちゃん\\
        & 「寄る」& ゆっず &ゆっじ & ゆっじゃん\\
        & 「切る」 &きー & きし & きしゃん\\
        & 「言う」 &いー & いじ & いじゃん \\
        & 「買う」 &こー & けー & かーん \\
        & 「結う」 &ゆー & いぇー & やーん \\
        %\hhline{~----}
        \rom{1}（変則） &「ある」& あい &－& ねーん \\
        & 「\ntilde{}だ」 &あい &－& あらん \\
        & 「居る」 &ぶい & ぶり & ぶら（ー）ん\\
        \hline
        \rom{2} &「見る」& みー & みーる & みーん\\
         & 「貰う」& いぇー & いぇーる & いぇーん \\
         & 「下る」& うり & うりる & うりん\\
        %\hhline{~---}
        \rom{2}（変則）&「寝る」&にに &ににる／にんる&ににん\\
        \midrule
        混合型 & 「死ぬ」&すに & すにる／すんる & すなん \\
        \midrule
        変則型 & 「する」&しー & しる & しゅ（ー）ん \\
         & 「来る」 &きー& くー  & く（ー）ん \\
     \bottomrule
 \end{tabular}
}
\end{table}

